\chapter{Solution}\label{ch:solution}

TODO

\section{Graph Schema}\label{sec:solution_schema}

\Cref{lst:cypher-example} shows an example Cypher query.

\begin{lstlisting}[language=Cypher,
                   caption={Example Cypher query (placeholder).},
                   label={lst:cypher-example}]
// Find failure modes linked to a component
MATCH (c:Component {name: $component})<-[:AFFECTS]-(f:FailureMode)
OPTIONAL MATCH (f)-[:CAUSED_BY]->(cause:Cause)
RETURN f.description AS failureMode,
       collect(DISTINCT cause.description) AS causes
ORDER BY failureMode
LIMIT 10
\end{lstlisting}

\section{Agent Architecture}\label{sec:solution_agent}

\Cref{lst:python-example} shows an example in Python.

\begin{lstlisting}[style=python,
                   caption={Example Python code (placeholder).},
                   label={lst:python-example}]
def run_query(driver, cypher: str, params: dict) -> list[dict]:
    """Execute a Cypher query and return the records."""
    with driver.session() as session:
        result = session.run(cypher, params)
        return [record.data() for record in result]
\end{lstlisting}

A prompt can be shown with the \texttt{prompt} style:

\begin{lstlisting}[style=prompt,
                   caption={Example prompt (placeholder).},
                   label={lst:prompt-example}]
You are an expert in Neo4j Cypher. Given the schema below,
translate the user's question into a single Cypher query.
\end{lstlisting}
